\documentclass[11pt]{article}
\usepackage{amsmath,amssymb,amsthm}
\usepackage{geometry}
\usepackage{hyperref}
\usepackage{enumitem}
\usepackage{graphicx}
\geometry{margin=1in}
\title{Connecting the Navier-Stokes Reflection Map to the Q_{NS} Parameter}
\author{}
\date{}
\begin{document}
\maketitle
\begin{abstract}
This document connects the Navier-Stokes reflection map to the Q_{NS} parameter, showing how scale-dependent quantities derived from the reflection map give rise to the Q_{NS} parameter that characterizes turbulence structure.
\end{abstract}

\section{1. From Reflection Map to Scale-Dependent Quantities}

Starting from the reflection map for enstrophy:
\[
\mathcal{R}[\mathcal{E}] = \underbrace{2(\mathbf{S} \cdot \boldsymbol{\omega}) \cdot \boldsymbol{\omega}}_{\text{Production}} - \underbrace{2\nu \|\nabla \boldsymbol{\omega}\|^2}_{\text{Dissipation}}
\]

We apply Littlewood-Paley decomposition to isolate scale-dependent components. Let \(\Delta_j\) be the Littlewood-Paley projector at scale \(j\) (wavelength \(\sim 2^j\)).

The scale-reflected and scale-inflicted components are:

\begin{itemize}
\item \textbf{Scale-Inflicted (Nonlinear Term)}:
  \[
  \mathcal{I}_j = \left\langle \Delta_j[(\mathbf{u} \cdot \nabla)\mathbf{u}] , \Delta_j\mathbf{u} \right\rangle
  \]
  This represents the energy flux through scale \(j\) from the nonlinear term.

\item \textbf{Scale-Reflected (Dissipative Term)}:
  \[
  \mathcal{R}_j = \nu \left\langle \Delta_j[-\Delta\mathbf{u}] , \Delta_j\mathbf{u} \right\rangle = \nu \|\nabla \Delta_j\mathbf{u}\|^2
  \]
  This represents the viscous dissipation at scale \(j\).
\end{itemize}

\section{2. Derivation of Q_{NS} from First Principles}

The energy flux through scale \(j\) can be related to our earlier definitions:
\[
\text{energy\_flux}(j) \sim \mathcal{I}_j
\]

The enstrophy at scale \(j\) relates to dissipation:
\[
\text{enstrophy\_LP}(j) \sim \|\nabla \Delta_j\mathbf{u}\|^2 \sim \frac{\mathcal{R}_j}{\nu}
\]

Therefore:
\[
Q_{NS}(j) = \frac{\text{energy\_flux}(j)}{\nu \cdot \text{enstrophy\_LP}(j) + 1} \approx \frac{\mathcal{I}_j}{\mathcal{R}_j + \nu}
\]

With appropriate normalization (setting \(\nu = 1\) in dimensionless form), we recover:
\[
Q_{NS}(j) \approx \frac{\mathcal{I}_j}{\mathcal{R}_j + 1}
\]

\section{3. The Reflection Map as Q_{NS} < 1 Condition}

The reflection map's effectiveness can be measured by:
\[
\text{Effectiveness}(j) = \frac{\mathcal{R}_j}{\mathcal{I}_j + \mathcal{R}_j}
\]

Then:
\[
Q_{NS}(j) < 1 \iff \mathcal{I}_j < \mathcal{R}_j + 1 \iff \text{Effectiveness}(j) > \frac{1}{2} \text{ (for large } \mathcal{I}_j, \mathcal{R}_j\text{)}
\]

More precisely, the reflection is effective (dissipation dominates) when:
\[
\mathcal{R}_j > \mathcal{I}_j \iff Q_{NS}(j) < \frac{\mathcal{R}_j}{\mathcal{R}_j + 1} < 1
\]

\section{4. Connection to Material Derivative Formulation}

Recall the material derivative of enstrophy:
\[
\frac{D\|\boldsymbol{\omega}\|^2}{Dt} = 2(\mathbf{S} \cdot \boldsymbol{\omega}) \cdot \boldsymbol{\omega} - 2\nu \|\nabla \boldsymbol{\omega}\|^2
\]

Integrating along fluid particle trajectories and applying Littlewood-Paley decomposition:
\[
\left\langle \Delta_j\left[\frac{D\|\boldsymbol{\omega}\|^2}{Dt}\right], 1 \right\rangle = \underbrace{2\left\langle \Delta_j[(\mathbf{S} \cdot \boldsymbol{\omega}) \cdot \boldsymbol{\omega}], 1 \right\rangle}_{\text{Scale-Inflicted}} - \underbrace{2\nu\left\langle \Delta_j[\|\nabla \boldsymbol{\omega}\|^2], 1 \right\rangle}_{\text{Scale-Reflected}}
\]

The scale-dependent Q_{NS} parameter emerges when we consider the ratio of these terms in the inertial range where the time derivative term becomes negligible.

\section{5. Implications for the Lean Formalization}

To make this rigorous in the NS.lean file, we should:

\begin{enumerate}
\item \textbf{Define the reflection map explicitly} as an operator on velocity fields
\item \textbf{Prove its relationship} to the material derivative of enstrophy (axiom NS_material_derivative_enstrophy)
\item \textbf{Show how scale decomposition} gives rise to the energy_flux and enstrophy_LP functions
\item \textbf{Establish that Q_{NS} < 1} corresponds to the reflection map being dissipative on average
\item \textbf{Replace the axiomatic treatment} with theorems derived from this reflection map framework
\end{enumerate}

This transforms the Q_{NS} parameter from an ad hoc definition to a direct consequence of the reflection map that characterizes turbulence structure.
\end{document}